\section{本章小结}
控制前先看执行器能否跟上参考轨迹。模型描述动态，容限描述可接受误差，代价函数权衡跟踪和控制开销，约束表达设备极限。MPC 每轮根据最新状态规划，只执行首个控制量。应熟练核对周期、单位、限幅与求解状态；求解器配置及实车稳定性需单独验证。

\begin{problemset}
\item 基础：y[k+1]=0.9y[k]+0.1u[k]，y[0]=0、u 恒为 1，前两步输出是多少？
\item 基础：纯延时 0.01 s 对 5 Hz 正弦信号的相位影响是多少？是否改变该信号幅值？
\item 基础：有效宽度 0.08 m、距离 4 m，小角度半角容限约为多少？
\item 基础：当前 y=0、参考 1，代价 (1-u)²+3u²，无约束最优 u 是多少？加入 
\item 基础：得到控制序列 [0.2,0.1,0.0] 后，滚动时域控制现在执行什么？下一周期做什么？
\item 提高：周期为 5 ms、预测步数 20，物理预测时域多长？求解 8 ms 会产生什么问题？
\item 实践：在本章滚动程序中第三轮更新后施加 -0.1 状态扰动，打印后续控制，确认每次都满足限幅。
\end{problemset}

\section{本章验收}
\begin{enumerate}
    \item 提交简化执行器的阶跃响应，标出响应时间与稳态误差，再演示改变模型参数后的变化。
    \item 运行受限滚动控制，对检查者指定的参考跳变与扰动展示状态、控制量和代价，解释跟踪项与控制项的作用。
    \item 模拟求解超时或不可行状态，展示日志和回退输出；全程只做离线仿真，不将教学模型直接用于开火。
\end{enumerate}
